% Block Puzzle RL — Beamer presentation
% Compile with: xelatex Block_Puzzle_RL.tex
\documentclass[aspectratio=169]{beamer}

\usepackage{fontspec}
\usepackage{polyglossia}
\setmainlanguage{russian}
\setotherlanguage{english}
\setmainfont{Inter}
\setsansfont{DejaVu Sans}
\setmonofont{DejaVu Sans Mono}
\usepackage{graphicx}
% multimedia несовместим с xelatex (\pdfmark), используем href
\usepackage{booktabs}
\usepackage{array}
\usepackage{xcolor}
\usepackage{colortbl}
\usepackage{tikz}
\usepackage{hyperref}

\graphicspath{{pictures/}}

% ── Color palette ─────────────────────────────────────────────────────────────
\definecolor{ThemeBG}{HTML}{0d0d1a}
\definecolor{ThemeCellFilled}{HTML}{4a90d9}
\definecolor{ThemeCellEmpty}{HTML}{1a1a30}
\definecolor{ThemeGridStroke}{HTML}{2a2a50}
\definecolor{ThemeGreen}{HTML}{27ae60}
\definecolor{ThemeBlue}{HTML}{2980b9}
\definecolor{ThemeRed}{HTML}{c0392b}
\definecolor{ThemeYellow}{HTML}{f39c12}
\definecolor{ThemePurple}{HTML}{9b59b6}
\definecolor{ThemeOrange}{HTML}{e67e22}
\definecolor{ThemeText}{HTML}{e8eaf6}
\definecolor{ThemeTextMuted}{HTML}{8899bb}

% ── Beamer chrome ─────────────────────────────────────────────────────────────
\usetheme{default}
\usecolortheme{default}
\setbeamertemplate{navigation symbols}{}

% Canvas & text
\setbeamercolor{background canvas}{bg=ThemeBG}
\setbeamercolor{normal text}{fg=ThemeText}
\setbeamercolor{alerted text}{fg=ThemeYellow}
\setbeamercolor{example text}{fg=ThemeGreen}

% Title page
\setbeamercolor{title}{fg=ThemeCellFilled}
\setbeamerfont{title}{size=\LARGE, series=\bfseries}
\setbeamercolor{subtitle}{fg=ThemeYellow}
\setbeamerfont{subtitle}{size=\normalsize}
\setbeamercolor{author}{fg=ThemeTextMuted}
\setbeamercolor{date}{fg=ThemeTextMuted}

% Frame title with separator rule
\setbeamercolor{frametitle}{fg=ThemeCellFilled, bg=ThemeBG}
\setbeamerfont{frametitle}{size=\large, series=\bfseries}
\setbeamertemplate{frametitle}{%
  \vspace{0.5ex}%
  \insertframetitle%
  \vspace{0.25ex}%
  \par{\color{ThemeGridStroke}\rule{\textwidth}{0.6pt}}%
  \vspace{-0.5ex}%
}

% Structure & bullets
\setbeamercolor{structure}{fg=ThemeCellFilled}
\setbeamercolor{itemize item}{fg=ThemeCellFilled}
\setbeamercolor{itemize subitem}{fg=ThemeGreen}
\setbeamercolor{itemize subsubitem}{fg=ThemeOrange}
\setbeamertemplate{itemize item}{\small\raise0.5pt\hbox{$\blacktriangleright$}}
\setbeamertemplate{itemize subitem}{\small\raise0.5pt\hbox{$\bullet$}}

% Blocks
\setbeamercolor{block title}{fg=ThemeCellFilled, bg=ThemeGridStroke}
\setbeamercolor{block body}{fg=ThemeText, bg=ThemeCellEmpty}

% Footline
\setbeamercolor{footline}{fg=ThemeTextMuted, bg=ThemeBG}
\setbeamerfont{footline}{size=\tiny}
\setbeamertemplate{footline}{%
  \hfill%
  {\tiny\color{ThemeTextMuted}\insertframenumber\,/\,\inserttotalframenumber}%
  \hspace{4mm}\vspace{2mm}%
}

% Table rules color
\arrayrulecolor{ThemeGridStroke}

% Hyperlink color
\hypersetup{colorlinks=true, urlcolor=ThemeCellFilled, linkcolor=ThemeCellFilled}

% ── Custom title page ─────────────────────────────────────────────────────────
\setbeamertemplate{title page}{%
	\begin{tikzpicture}[remember picture, overlay]
		
		% ── Clean background ──────────────────────────────────────────────────
		\fill[ThemeBG] (current page.south west)
		rectangle (current page.north east);
		
		% ── Subtle accent line ────────────────────────────────────────────────
		\fill[ThemeBlue!60]
		([yshift=-2cm]current page.north west)
		rectangle ([yshift=-1.8cm]current page.north east);
		
	\end{tikzpicture}%
	
	% ── Top spacing (уменьшили, чтобы поднять контент) ──────────────────────
	\vspace{1.5cm}
	
	\begin{flushleft}
		
		% ── Title ─────────────────────────────────────────────────────────────
		{\usebeamerfont{title}
			\usebeamercolor[fg]{title}
			\fontsize{34}{38}\selectfont
			\inserttitle\par}
		
		\vspace{0.8em}
		
		% ── Subtitle ──────────────────────────────────────────────────────────
		{\usebeamerfont{subtitle}
			\color{ThemeGridStroke}
			\large
			\insertsubtitle\par}
		
	\end{flushleft}
	
	% ── Заполняем пространство до низа ──────────────────────────────────────
	\vfill
	
	% ── Bottom block (автор + год) ──────────────────────────────────────────
	\begin{flushleft}
		{\small \color{ThemeGridStroke}
			Автор: Андрей Сергиенко Б-82{\hspace{4cm}} \\
			Год: 2026{\hspace{2cm}}}
	\end{flushleft}
	
	\vspace{0.8cm}
}

\title{Block Puzzle: Обучение с подкреплением}
\subtitle{MaskablePPO · 3 поколения агентов · MLP + SmallCNN}
\date{2026}

\begin{document}

% ──────────────────────────────────────────────────────────────────────────────
\begin{frame}
  \titlepage
\end{frame}

% ──────────────────────────────────────────────────────────────────────────────
\begin{frame}{Постановка задачи}
	\begin{columns}
		\column{0.45\textwidth}
		\centering
		\includegraphics[width=0.8\textwidth,keepaspectratio]{board_game_demo}

		\column{0.55\textwidth}
		\vspace{-0.5cm}
		\begin{itemize}
			\item Поле 8×8, тройка случайных фигур, заполненные строки/столбцы очищаются

			\item \textbf{Пространство действий:} до $3 \times 8 \times 8 = 192$, часть невалидна

			\item \textbf{Разреженная награда:} линии очищаются редко

			\item \textbf{Планирование:} решения внутри хода зависят друг от друга

			\item \textbf{Стохастика:} следующая тройка фигур случайна
		\end{itemize}
	\end{columns}
\end{frame}

% ──────────────────────────────────────────────────────────────────────────────
\begin{frame}{Обучение с подкреплением: формальная постановка}
  \begin{columns}[t]
    \begin{column}{0.50\textwidth}
      \textbf{Марковский процесс принятия решений (MDP)}

      \medskip
      $\langle \mathcal{S},\, \mathcal{A},\, \mathcal{T},\, R,\, \gamma \rangle$

      \medskip
      \begin{itemize}
        \item $s_t \in \mathcal{S}$ — состояние: поле $8{\times}8$ + фигуры
        \item $a_t \in \mathcal{A}(s_t)$ — \emph{валидное} действие
        \item $r_t = R(s_t, a_t)$ — мгновенная награда
        \item $\gamma = 0.99$ — дисконт будущих наград
      \end{itemize}

      \medskip
      \textbf{Цель:} максимизировать ожидаемый доход
      \[
        G_t = \sum_{k=0}^{\infty} \gamma^k\, r_{t+k}
      \]
    \end{column}
    \begin{column}{0.46\textwidth}
      \textbf{Агент и среда}

      \medskip
      \begin{itemize}
        \item Агент наблюдает $s_t$, выбирает действие $a_t$
        \item Среда переходит в $s_{t+1}$ и возвращает $r_t$
        \item Агент обновляет политику $\pi_\theta$, чтобы максимизировать $G_t$
        \item Цикл повторяется до конца эпизода
      \end{itemize}

      \medskip
      \textbf{Ключевая сложность:}\\
      агент должен самостоятельно понять, какие действия ведут к увеличению награды
    \end{column}
  \end{columns}
\end{frame}

% ──────────────────────────────────────────────────────────────────────────────
\begin{frame}{Обучение с подкреплением: PPO + Action Masking}
  \begin{columns}[t]
    \begin{column}{0.50\textwidth}
      \textbf{Policy Gradient / PPO}

      \medskip
      Политика $\pi_\theta(a \mid s)$ — нейросеть с весами $\theta$.

      \medskip
      PPO ограничивает шаг через \emph{clipped surrogate}:
      \[
        L^{\text{CLIP}} = \mathbb{E}\!\left[\min\!\left(
          \rho_t \hat{A}_t,\;
          \mathrm{clip}(\rho_t, 1{\pm}\varepsilon)\,\hat{A}_t
        \right)\right]
      \]
      $\rho_t = \dfrac{\pi_\theta(a_t|s_t)}{\pi_{\theta_\text{old}}(a_t|s_t)}$,\quad $\varepsilon = 0.2$\\
      \medskip
      Return: $G_t = \sum_{k=0}^{\infty} \gamma^k r_{t+k}$
    \end{column}
    \begin{column}{0.46\textwidth}
      \textbf{Action masking}

      \medskip
      Невалидные действия исключаются явно — не штрафом, а обнулением вероятности:
      \[
        \tilde{\pi}_\theta(a \mid s) = \frac{\pi_\theta(a|s)\cdot m(a,s)}
          {\sum_{a'} \pi_\theta(a'|s)\cdot m(a',s)}
      \]
      $m(a,s) \in \{0,1\}$ — маска допустимых ходов

      \medskip
      Это ускоряет обучение и исключает пустые градиенты от заведомо невозможных ходов.
    \end{column}
  \end{columns}
\end{frame}

% ──────────────────────────────────────────────────────────────────────────────
\begin{frame}{Алгоритм: MaskablePPO}
  \begin{itemize}
    \item \textbf{PPO} (Proximal Policy Optimization) — один из самых надёжных on-policy алгоритмов
    \item \textbf{MaskablePPO} из sb3-contrib — невалидные действия исключаются из вероятностного распределения явно, не через штраф
  \end{itemize}

  \bigskip
  \textbf{Ключевые метрики обучения}
  \begin{itemize}
    \item \textbf{explained\_variance (EV):} >0.3 — обучение идёт верно, 1.0 — идеал, <0 — хуже случайного
    \item \textbf{clip\_fraction:} норма 0.03–0.08; выше — политика меняется слишком резко
    \item \textbf{entropy\_loss:} слишком низкая — агент «зациклился» на одних действиях
  \end{itemize}
\end{frame}

% ──────────────────────────────────────────────────────────────────────────────
\begin{frame}{Гиперпараметры}
  \begin{center}
  \begin{tabular}{@{}lll@{}}
    \toprule
    \rowcolor{ThemeGridStroke}
    \color{ThemeCellFilled}Параметр & \color{ThemeCellFilled}Значение & \color{ThemeCellFilled}Описание \\
    \midrule
    n\_steps        & 512   & шагов перед обновлением политики \\
    n\_epochs       & 4     & проходов по собранным данным \\
    learning\_rate  & 1e-4  & скорость обучения (MLP); 7e-5 для CNN \\
    n\_envs         & 16    & параллельных сред \\
    total\_steps    & 5 M   & шагов обучения всего \\
    batch\_size     & 256   & размер минибатча \\
    ent\_coef       & 0.005 & коэффициент энтропийного бонуса \\
    \bottomrule
  \end{tabular}
  \end{center}

  \bigskip
  \textbf{Hyperparameter sweep:} 7 конфигураций × 100k шагов, оценка на 100 эпизодах. Были найдены оптимальные параметры
\end{frame}

% ──────────────────────────────────────────────────────────────────────────────
\begin{frame}{Окружение: Gymnasium + наблюдение}
  \begin{columns}[t]
    \begin{column}{0.48\textwidth}
      \textbf{Gymnasium — стандарт RL-окружений}

      \medskip
      \begin{itemize}
        \item Интерфейс: \texttt{reset()} / \texttt{step(action)}
        \item \texttt{step} возвращает $(s', r, \text{done}, \text{trunc}, \text{info})$
        \item \texttt{SubprocVecEnv} — 16 сред в параллельных процессах
        \item \texttt{ActionMasker} (sb3-contrib) — передаёт маску допустимых ходов в MaskablePPO
      \end{itemize}

      \medskip
      \textbf{Игра → нейросеть: тензор $(C, 8, 8)$}
    \end{column}
    \begin{column}{0.48\textwidth}
      \textbf{Почему тензор, а не вектор?}

      \medskip
      Плоский вектор теряет пространственную структуру: \\
      клетка $(3,5)$ и $(3,6)$ оказываются далеко друг от друга в памяти.

      \medskip
      Тензор $\mathbb{R}^{C \times 8 \times 8}$ сохраняет соседство —
      свёрточные фильтры CNN видят локальные паттерны (уголки, линии, дыры).
    \end{column}
  \end{columns}
\end{frame}

% ──────────────────────────────────────────────────────────────────────────────
\begin{frame}{Инженерия наград}
  \begin{columns}[t]
    \begin{column}{0.50\textwidth}
      \textbf{Компоненты функции награды}

      \bigskip
      \begin{tabular}{@{}lll@{}}
        \toprule
        \rowcolor{ThemeGridStroke}
        \color{ThemeCellFilled}Компонент & \color{ThemeCellFilled}Значение & \\
        \midrule
        place\_piece    & +0.2 \\
        cell\_potential & +varies \\
        lines\_cleared  & +1.0×N  \\
        game\_over      & −1.0 \\
        round\_complete & +0.5 (Gen 3)\\
        \bottomrule
      \end{tabular}
    \end{column}
    \begin{column}{0.46\textwidth}
      \textbf{Ключевая идея: cell\_potential}

      \smallskip
      \[
      r \mathrel{+}= k \cdot \frac{\sum (row\_fill + col\_fill)}{2}
      \]

      \bigskip
      Агент самостоятельно учится ставить фигуры в почти заполненные строки.
      Каналы 4–5 показывают заполненность,
      \texttt{cell\_potential} вознаграждает за её использование.
    \end{column}
  \end{columns}
\end{frame}

% ──────────────────────────────────────────────────────────────────────────────
\begin{frame}{Бейзлайны}
  \begin{columns}[t]
    \begin{column}{0.50\textwidth}
      \textbf{RandomAgent} \\
	  \medskip
      Выбирает случайный валидный ход из маски действий.\\
      \bigskip
      Lines: 1.79 ± 2.12 \quad Reward: 10.48 ± 11.67 \quad CV = 1.18
    \end{column}
    \begin{column}{0.46\textwidth}
      \textbf{HeuristicAgent} (Dellacherie, Tetris 2003)\\
      \medskip
      На каждом ходу симулирует все валидные размещения,
      оценивает доску по 7 признакам:
      очищенные линии, дыры (BFS от границ),
      почти заполненные линии, компактность, дисперсия заполнения.\\
      \bigskip
      Lines: 8.48 ± 7.60 \quad Reward: 45.91 ± 40.10 \quad CV = 0.90
    \end{column}
  \end{columns}
\end{frame}

% ──────────────────────────────────────────────────────────────────────────────
\begin{frame}{SmallCNN: архитектура}
  NatureCNN (встроенная в SB3) использует stride=4/2 — на поле 8×8 схлопывается в ошибку.

  \bigskip
  \begin{center}
  \begin{tabular}{@{}lll@{}}
    \toprule
    \rowcolor{ThemeGridStroke}
    \color{ThemeCellFilled}Слой & \color{ThemeCellFilled}Выход & \color{ThemeCellFilled}Примечание \\
    \midrule
    Вход                     & [B, C, 8, 8]  & C каналов \\
    Conv2d(C→32, 3×3, s=1)   & [B, 32, 8, 8] & stride=1, padding=1 \\
    Conv2d(32→64, 3×3, s=1)  & [B, 64, 8, 8] & \\
    Conv2d(64→64, 3×3, s=1)  & [B, 64, 8, 8] & рецептивное поле 7×7 \\
    Flatten                  & [B, 4096]     & \\
    Linear(4096→256) + ReLU  & [B, 256]      & размер выводится авто \\
    MLP actor / critic       & —             & поверх 256-d эмбеддинга \\
    \bottomrule
  \end{tabular}
  \end{center}
\end{frame}

% ──────────────────────────────────────────────────────────────────────────────
\begin{frame}{SmallCNN: ключевые решения}
  \begin{itemize}
    \item \textbf{stride=1, padding=1:} размер 8×8 сохраняется через все свёрточные слои; рецептивное поле 7×7 — почти вся доска

    \item \textbf{BatchNorm убран:} \texttt{eval()} vs \texttt{train()} — разные режимы BN искажают advantage-оценки.
      Результат с BN: clip\_fraction=0.45, EV=−0.57 (хуже случайного)

    \item \textbf{Авто-размер Linear:} \texttt{flat\_size = cnn(zeros(1,C,8,8)).shape[1]} — инвариантен к количеству каналов и размеру доски

    \item \textbf{Learning rate критичен:}
      lr=3e-5 → clip\_fraction=0.002 (агент не учится);
      lr=7e-5 → нормальная сходимость
  \end{itemize}
\end{frame}

\begin{frame}{SmallCNN: демонстрация}
  \begin{center}
    \href{run:video/CNNAnatomyScene.mp4}{%
      \includegraphics[width=0.85\textwidth,height=0.75\textheight,keepaspectratio]{cnn_preview}%
    }
  \end{center}
\end{frame}

% ──────────────────────────────────────────────────────────────────────────────
\begin{frame}{Поколение 1: базовая архитектура}
  Наблюдение: \textbf{(6, 8, 8)} — поле + 3 формы фигур + заполненность строк/столбцов

  \bigskip
  \begin{itemize}
    \item \textbf{MLP:} тензор → вектор 384 → actor [256, 256] / critic [512, 512]
    \item \textbf{CNN:} SmallCNN (256-d эмбеддинг) → те же MLP-головы
  \end{itemize}

	\begin{tabular}{@{}lrrrrr@{}}
		\toprule
		\rowcolor{ThemeGridStroke}
		\color{ThemeCellFilled}Метрика & \color{ThemeCellFilled}Random & \color{ThemeCellFilled}Heuristic & \color{ThemeCellFilled}MLP Gen 1 & \color{ThemeCellFilled}CNN Gen 1 \\
		\midrule
		Lines mean  & 1.79   & 8.48      & \textbf{8.84}  & 7.36 \\
		Lines std   & 2.12   & 7.60      & 5.90           & \textbf{4.74} \\
		Lines CV    & 1.18   & 0.90      & \textbf{0.67}  & \textbf{0.64} \\
		Reward mean & 10.48  & 45.91     & \textbf{48.23} & 39.92 \\
		Reward std  & 11.67  & 40.10     & 31.45          & \textbf{25.23} \\
		\bottomrule
	\end{tabular}

	\bigskip
	MLP превзошёл эвристику: +4.2\% по lines.
	Оба нейросетевых агента стабильнее: CV 0.64–0.67 vs 0.90 (−28\%).
\end{frame}


% ──────────────────────────────────────────────────────────────────────────────
\begin{frame}{Поколение 2: Placement Heatmap}
  Наблюдение: \textbf{(9, 8, 8)} — добавлены каналы 6–8: placement heatmap для каждой фигуры.

  \bigskip
  Ячейка (y,x) = 1.0, если фигуру можно поставить с верхним левым углом в (y,x).
  Overhead нулевой — маска action\_mask всё равно считается для MaskablePPO.

  \bigskip
  Gen 1 → декларативно: «вот форма фигуры» \\
  Gen 2 → процедурно: «вот куда её можно поставить»

  \bigskip
  \begin{tabular}{@{}lrrrrr@{}}
    \toprule
    \rowcolor{ThemeGridStroke}
    \color{ThemeCellFilled}Метрика & \color{ThemeCellFilled}Random & \color{ThemeCellFilled}Heuristic & \color{ThemeCellFilled}MLP Gen 2 & \color{ThemeCellFilled}CNN (lr=7e-5) \\
    \midrule
    Lines mean   & 1.79   & 8.48      & \textbf{9.88}   & 8.50 \\
    Lines CV     & 1.18   & 0.90      & \textbf{0.67}            & \textbf{0.67} \\
    Reward mean  & 10.48  & 45.91     & \textbf{53.89}  & 46.36 \\
    Pieces placed & 16.32 & 32.31     & \textbf{35.85}  & 32.68 \\
    \bottomrule
  \end{tabular}
\end{frame}

% ──────────────────────────────────────────────────────────────────────────────
\begin{frame}{Поколение 3: Топологическое зрение}
  Наблюдение: \textbf{(14, 8, 8)} — добавлены каналы 9–13.

  \bigskip
  Проблема Gen 2: агент не различал «30 пустых клеток = живая доска» от
  «30 пустых клеток = 30 изолированных фрагмента».

  \bigskip
  \begin{itemize}
    \item \textbf{Каналы 9–11 — Survivability:}
      для каждой валидной позиции фигуры i — сколько ходов остаётся у двух других фигур после этого размещения.
      Нормировано [0,1]. Lookahead внутри раунда.

    \item \textbf{Канал 12 — Largest Blob:}
      BFS по пустым клеткам. Крупнейшая связная компонента = 1.0, остальные клетки = 0.

    \item \textbf{Канал 13 — Dead Zone:}
      BFS-компонента мертва, если ни одна из текущих трёх фигур не может её затронуть.
      Проверяем текущий пул фигур.
  \end{itemize}

  \bigskip
  \textbf{round\_complete = +0.5} — явный сигнал «дожить до следующей тройки = хорошо»
\end{frame}

% ──────────────────────────────────────────────────────────────────────────────
\begin{frame}{Поколение 3: Transfer Learning для CNN}
  CNN не обучается с нуля на 14 каналах:
  Conv1 применяет фильтры ко всем 14 каналам одновременно —
  шум из topology-каналов портит фильтры базовой геометрии.
  CNN с нуля: \textbf{4.73 lines} (Gen 2 CNN: 8.50).

  \bigskip
  \textbf{Решение — Transfer Learning:}
  \begin{itemize}
    \item Conv1: (32, 9, 3×3) → (32, 14, 3×3)
    \item каналы 0–8 = веса из Gen 2; каналы 9–13 = нули
    \item Conv2, Conv3, Linear, MLP-головы — копируются без изменений
  \end{itemize}

  \bigskip
  \begin{tabular}{@{}lrrrr@{}}
    \toprule
    \rowcolor{ThemeGridStroke}
    \color{ThemeCellFilled}Метрика & \color{ThemeCellFilled}CNN Gen 2 & \color{ThemeCellFilled}CNN Gen 3 (с нуля) & \color{ThemeCellFilled}CNN Gen 3 (TL) \\
    \midrule
    Lines mean  & 8.50      & 4.73               & \textbf{8.85} \\
    Lines std   & 5.72      & 3.79               & \textbf{5.66} \\
    Reward mean & 46.36     & 26.63              & \textbf{48.23} \\
    Pieces      & 32.68     & 23.63              & \textbf{33.52} \\
    \bottomrule
  \end{tabular}
\end{frame}

\begin{frame}{Что по итогу видит агент?}
	\begin{center}
		\href{run:video/ObsEvolutionScene.mp4}{%
			\includegraphics[width=0.85\textwidth,height=0.75\textheight,keepaspectratio]{cnn_preview}%
		}
	\end{center}
\end{frame}

% ──────────────────────────────────────────────────────────────────────────────
\begin{frame}{Итоговые результаты MLP}
  \begin{tabular}{@{}lrrrrrr@{}}
    \toprule
    \rowcolor{ThemeGridStroke}
    \color{ThemeCellFilled}Метрика & \color{ThemeCellFilled}Random & \color{ThemeCellFilled}Heuristic & \color{ThemeCellFilled}MLP Gen 1 & \color{ThemeCellFilled}MLP Gen 2 & \color{ThemeCellFilled}MLP Gen 3 \\
    \midrule
    Lines mean    & 1.79   & 8.48      & 8.84      & 9.88      & \textbf{11.21} \\
    Lines std     & 2.12   & 7.60      & 5.90      & 6.57      & 7.68 \\
    Lines CV      & 1.18   & 0.90      & 0.67      & 0.67      & 0.68 \\
    Reward mean   & 10.48  & 45.91     & 48.23     & 53.89     & \textbf{60.76} \\
    Reward std    & 11.67  & 40.10     & 31.45     & 35.13     & 40.96 \\
    Pieces placed & 16.32  & 32.31     & 33.43     & 35.85     & \textbf{38.81} \\
    \bottomrule
  \end{tabular}

  \bigskip
  MLP Gen 3 vs Heuristic: \textbf{+32\% по lines cleared}, \textbf{+32\% по reward}
\end{frame}

% ──────────────────────────────────────────────────────────────────────────────
\begin{frame}{Кривые обучения MLP}
	\begin{center}
		\includegraphics[height=0.80\textheight,keepaspectratio]{chart_learning_mlp}
	\end{center}
\end{frame}

% ──────────────────────────────────────────────────────────────────────────────
\begin{frame}{Итоговые результаты CNN}
  \footnotesize
  \begin{tabular}{@{}lrrrrrr@{}}
    \toprule
    \rowcolor{ThemeGridStroke}
    \color{ThemeCellFilled}Метрика & \color{ThemeCellFilled}Random & \color{ThemeCellFilled}Heuristic & \color{ThemeCellFilled}CNN Gen 1 & \color{ThemeCellFilled}CNN Gen 2 & \color{ThemeCellFilled}CNN Gen 3 TL \\
    \midrule
    Lines mean     & 1.79   & 8.48      & 7.36      & 8.50            & \textbf{8.85} \\
    Lines std      & 2.12   & 7.60      & 4.74      & 5.72               & \textbf{5.66} \\
    Reward mean    & 10.48  & 45.91     & 39.92     & 46.36             & \textbf{48.23} \\
    Reward std     & 11.67  & 40.10     & 25.23     & 30.50             & \textbf{30.17} \\
    Pieces placed  & 16.32  & 32.31     & 29.40     & 32.68             & \textbf{33.52} \\
    \bottomrule
  \end{tabular}

  \bigskip
  \begin{itemize}
    \item CNN Gen 3 \textbf{с нуля}: катастрофа — 4.73 lines (хуже Gen\,1)
    \item CNN Gen 3 \textbf{Transfer Learning}: восстанавливает уровень Gen\,2 и чуть превосходит его
    \item Потолок CNN на 8×8 — уровень эвристики; MLP Gen 3 уходит на \textbf{+32\%}
  \end{itemize}
\end{frame}

% ──────────────────────────────────────────────────────────────────────────────
\begin{frame}{Кривые обучения CNN}
	\begin{center}
		\includegraphics[height=0.80\textheight,keepaspectratio]{chart_learning_cnn}
	\end{center}
\end{frame}

% ──────────────────────────────────────────────────────────────────────────────
\begin{frame}{MLP vs CNN}
	\begin{center}
		\includegraphics[height=0.82\textheight,keepaspectratio]{chart_mlp_vs_cnn}
	\end{center}
\end{frame}

% ──────────────────────────────────────────────────────────────────────────────
\begin{frame}{Статистическая значимость}
  1000 независимых эпизодов · seed=42 · \textbf{MLP Gen 3 vs Heuristic}

  \bigskip
  \begin{itemize}
    \item \textbf{Mann-Whitney U: p < 0.001} — превосходство MLP значимо при любом уровне значимости
    \item \textbf{CLES = 0.634} — агент выигрывает в 63.4\% случайно выбранных пар эпизодов
    \item \textbf{Cohen's d = +0.348} — умеренный эффект (d > 0.2 считается малым)
    \item \textbf{Paired test (одинаковые фигуры): win rate 59–61\%} — медианная разница +2 линии и +13 очков
  \end{itemize}

  \bigskip
  \begin{tabular}{@{}lrrrrrr@{}}
    \toprule
    \rowcolor{ThemeGridStroke}
    \color{ThemeCellFilled}Агент & \color{ThemeCellFilled}≥1 & \color{ThemeCellFilled}≥5 & \color{ThemeCellFilled}≥10 & \color{ThemeCellFilled}≥15 & \color{ThemeCellFilled}≥20 & \color{ThemeCellFilled}≥30 \\
    \midrule
    Random        & 66\% & 12\% & 2\%  & 0\%  & 0\% & 0\% \\
    Heuristic     & 94\% & 64\% & 33\% & 17\% & 9\% & 3\% \\
    \textbf{MLP Gen 3} & \textbf{99\%} & \textbf{85\%} & \textbf{49\%} & \textbf{26\%} & \textbf{14\%} & 3\% \\
    CNN Gen 3 TL  & 99\% & 78\% & 40\% & 15\% & 6\% & 1\% \\
    \bottomrule
  \end{tabular}
\end{frame}

% ──────────────────────────────────────────────────────────────────────────────
\begin{frame}{Выводы}
  \begin{itemize}
    \item \textbf{Нейросеть превзошла вручную оптимизированный алгоритм} без единой строки знаний о правилах игры.
      Эвристика: 7 признаков + подобранные веса + BFS.
      Нейросеть: только состояние и награда.

    \item \textbf{Стабильность принципиальна:} CNN Gen 3: CV ≈ 0.64 vs \textasciitilde0.90 у эвристики (−29\%).
      Эвристика нестабильна по дизайну — жадный одношаговый выбор = лотерея при сложных раскладах.

    \item \textbf{Масштабируемость:} сложность эвристики на доске 16×16 растёт в 16×.
      Нейросеть: фиксированный граф вычислений — та же задержка.

    \item \textbf{Инференс в 10× быстрее:} Heuristic 77 сек vs MLP 7 сек на 500 эпизодов.

    \item \textbf{Feature engineering в RL работает:}
      Gen 1→2: +12\% по lines.
      Gen 2→3: +13\% по lines.
      Дать агенту готовую информацию эффективнее, чем ждать пока он выведет её сам.
  \end{itemize}
\end{frame}

% ──────────────────────────────────────────────────────────────────────────────
\begin{frame}{Почему CNN хуже MLP на 8×8?}
  \begin{columns}[t]
    \begin{column}{0.50\textwidth}
      \textbf{Корень проблемы — слишком маленькое поле}\\
      \smallskip
      На поле 8×8 уже три слоя Conv\,3×3 покрывают рецептивное поле 7×7 —
      \textbf{почти всю доску целиком}. CNN фактически вынуждена делать то же,
      что и MLP, но с большим числом параметров и overhead'ом свёрток.
      MLP при этом видит весь вектор сразу, без индуктивного смещения
      на локальность — и выигрывает именно там, где нужна \emph{глобальная} картина.
    \end{column}
    \begin{column}{0.46\textwidth}
      {\color{ThemeText}\textbf{Что изменится на 16×16?}}

      \medskip
      \begin{itemize}
        \item Рецептивное поле 7×7 покрывает лишь 19\% доски — CNN снова имеет смысл
        \item Пространство действий вырастет до $3 \times 16 \times 16 = 768$
        \item MLP вектор: $C \times 256 = 3584$ — размерность взрывается линейно
        \item CNN эмбеддинг остаётся 256-d вне зависимости от размера доски
        \item Эвристика замедляется в $16\times$, CNN — нет
      \end{itemize}
    \end{column}
  \end{columns}
\end{frame}

% ──────────────────────────────────────────────────────────────────────────────
\begin{frame}[shrink=8]{Что дальше?}
  \begin{columns}[t]
    \begin{column}{0.50\textwidth}
      {\color{ThemeText}\textbf{Скорость: среда на C + CUDA}}

      \smallskip
      \begin{itemize}
        \item \textbf{16 сред на Python — мало и медленно:}
          GIL, интерпретатор, \texttt{SubprocVecEnv} = 16 процессов с IPC-overhead'ом
        \item \textbf{Среда на C (CUDA-ядра):} тысячи сред живут прямо
          в видеопамяти — \texttt{step()} выполняется батчем на GPU без
          копирования данных между CPU и GPU
        \item \textbf{IsaacGym / EnvPool} — готовые фреймворки для GPU-сред:
          10\,000+ параллельных сред вместо 16
      \end{itemize}
    \end{column}
    \begin{column}{0.46\textwidth}
      {\color{ThemeText}\textbf{Архитектура и алгоритмы}}

      \medskip
      \begin{itemize}
        \item \textbf{Доска 16×16 / 10×10:} проверить гипотезу о CNN на большом поле
        \item \textbf{DQN/ Transformer:} попробовать другие архитектуры
        \item \textbf{Curriculum learning:} начинать с простых раскладов фигур,
          постепенно усложнять — ускоряет начальную сходимость
      \end{itemize}
    \end{column}
  \end{columns}
\end{frame}

\end{document}
